\subsection{Four cross states opposite a right-zero corner}

\begin{lemma}[The cofiber of a connected graph correspondence]
\label{C:connected-graph-cone-degree}
Let a finite monoid $M$ act on the right on a connected finite
graph $\Gamma$ by cellular maps, and suppose $H_1(\Gamma;\mathbb Q)$
is one-dimensional.  Write $\rho_m$ for the induced degree of
$m$ on that vector space.  Put
\[
 X=\operatorname{hocofib}(\Gamma_{hM}\longrightarrow BM).
\]
Then
\[
 \widetilde H_n(X;\mathbb Q)
   \cong\operatorname{Tor}^{\mathbb QM}_{n-2}
       (H_1(\Gamma;\mathbb Q),\mathbb Q)\qquad(n\geq2).
\]
In particular $H_2(X;\mathbb Q)\ne0$ forces $\rho_m=1$ for
every $m$.  If an $M$-equivariant self-map of $\Gamma$ has
degree $d\notin\{0,1\}$, then every positive-degree rational
cohomology product of $X$ is zero.
\end{lemma}
\begin{proof}
The augmented cellular chain complex
$C_1(\Gamma)\to C_0(\Gamma)\to\mathbb Q$ is a complex of
right $\mathbb QM$-modules.  Its only nonzero homology is
$H_1(\Gamma)$ in the first position.  Inserting this kernel is
an equivariant quasi-isomorphism.  Tensor with a free left
bar resolution and take the mapping cone of augmentation.
This computes the indicated homotopy cofiber and shifts the
kernel by two, proving the displayed formula.  It also gives
$\widetilde H_1(X;\mathbb Q)=0$.
In degree two the result is the coinvariant space
$H_1(\Gamma)\otimes_{\mathbb QM}\mathbb Q$; a single degree
$\rho_m\ne1$ makes it zero.

An equivariant degree-$d$ self-map acts as multiplication by
$d$ on the kernel in this chain comparison.  It therefore acts
as multiplication by $d$ on every positive homology and
cohomology group of the cofiber.  For positive-degree classes
$u,v$, naturality and multiplicativity give
$d(u v)=d^2(u v)$.  Over $\mathbb Q$ this forces $uv=0$.
\end{proof}

\begin{theorem}[Four-state right-zero reduction]
\label{C:right-zero-four-state-reduction}
Let $R=\{1,p,q\}$ with $ab=b$ for $a,b\in\{p,q\}$, and
let $S$ be a nonempty $(R,M)$-biset with $|S|\leq4$.
If its two-object collage has type $\mathbb{CP}^2$, an ordinary
right $M$-act cone on at most four states has that type.
\end{theorem}
\begin{proof}
The cases with at most three states follow from
Lemma~\ref{C:three-state-right-zero-biset}, the common-fixed-state
argument, and the finite-line lemma.  Suppose $|S|=4$ and let
$P,Q$ be the two retractions onto their common image $U$.
Use the natural graph model of
Lemma~\ref{C:right-zero-action-graph}: its vertices are $U$ and
its edges are the states outside $U$, with edge $s$ directed
from $Ps$ to $Qs$.  The graph carries the commuting right
$M$-action.  Since $BR$ is contractible, the collage is the
homotopy cofiber of $\Gamma_{hM}\to BM$.

If $|U|=1$, its point is fixed by both monoids, and every
degree-two product on the collage vanishes.  If $|U|=4$, the
left action is trivial, and the collage is equivalent to the
ordinary four-state $M$-act cone.
If $|U|=3$ and its single graph edge is not a loop, each graph
component is contractible.  The equivariant map to its two-point
component set is a homotopy equivalence, also after taking the
action bar.  This gives an ordinary two-state act cone.

Whenever $P=Q$, including the remaining $|U|=3$ case, the map
$P:S\to U$ is a biset retraction.  Let $Y$ be the collage on
$U$.  It is equivalent to the ordinary $M$-act cone on $U$,
and is a homotopy retract of the original collage $X$.
The relative quotient is $\Sigma(E_S/E_U)$, as in
Lemma~\ref{C:three-state-right-zero-biset}.
If $X\simeq\mathbb{CP}^2$ and $H_2(Y;\mathbb Z)=0$, every
rational degree-two class of $X$ is relative, so its square
would vanish.  Otherwise split integral homology makes
$H_2(Y)\to H_2(X)$ an isomorphism.  The nonzero primitive square
then makes $H_4(Y)\to H_4(X)$ an isomorphism as well.
All other reduced homology groups vanish, and $\pi_1(Y)=1$
because it is a retract of $\pi_1(X)=1$.  Homological Whitehead
makes $Y\to X$ a homotopy equivalence.  Thus this case gives
the required ordinary act cone.

It remains that $|U|=2$ and the two-edge graph is connected.
It has one-dimensional first homology.  A hypothetical
projective plane forces every actor to have degree $+1$ by
Lemma~\ref{C:connected-graph-cone-degree}.
There are three possible graph types.

If there is one connecting edge and one loop, a degree-one
actor must carry the unique loop to itself.  It therefore fixes
its vertex, which is also fixed by $R$.  This common fixed state
makes all degree-two classes relative to the two vertex spaces,
and their products vanish.  The connecting edge need not be
invariant: it may map to the loop.  Only the fixed loop vertex
is used.

If the two edges connect the two vertices in the same direction,
a degree-one actor must fix both vertices and both edges.
Swapping the edges has degree $-1$, while folding or collapsing
them has degree zero.  Hence there is again a common fixed
state and zero degree-two products.

Finally suppose the two connecting edges have opposite directions.
A degree-one actor is either the identity or the half-rotation
interchanging the two vertices and edges.  Identify this graph
with the circle so that the half-rotation is $z\mapsto-z$.
The degree-three map $z\mapsto z^3$ commutes with every actor.
Lemma~\ref{C:connected-graph-cone-degree} makes every positive
rational product zero, again excluding a projective plane.
These alternatives exhaust the four-state graph and finish
the reduction.
\end{proof}

\begin{corollary}\label{C:right-zero-four-state-bounds}
If $|M|\leq12$, a biset collage with the right-zero corner $R$
and at most four cross arrows cannot have type $\mathbb{CP}^2$.
For an arbitrary three-element opposite corner, the same
conclusion holds when $|M|\leq10$ and there are four cross arrows.
In particular the sixteen-arrow profile $(9,3,4)$ is excluded
without a new enumeration.
\end{corollary}
\begin{proof}
For $R$, use the reduction and
Corollary~\ref{C:four-state-cones-seventeen},
Theorem~\ref{C:three-state-cone-minimum22}, and the one- and
two-state finite-line cases.
Every other three-element monoid has projective trivial right
rational module, as checked by the splitting vectors in the
seven-monoid catalog of the preceding biset calculation.  The actual
orbit-act functor is then a rational homology and ring
equivalence.  For $|M|\leq9$ apply
Lemma~\ref{C:nine-bounded-rank-one}; for $|M|=10$ apply
Lemma~\ref{C:ten-rational-act-cones}.  Both use precisely the
bounded rational homology and one-dimensional degree-two
cohomology of a hypothetical projective plane.
\end{proof}
